\unit{Beam splitters, HOM interference and detectors}
\unitref{Weeks 2--3; optics tutorials; 2024--25 exams}
\subsection{How an interferometer mixes amplitudes}
A lossless two-port beam splitter transforms field
\emph{operators}, not photon probabilities:
\[
\boxed{
\begin{pmatrix}c\\d\end{pmatrix}=
\begin{pmatrix}t&r\\-r&t\end{pmatrix}
\begin{pmatrix}a\\b\end{pmatrix},
\quad t^2+r^2=1.}
\]
The numbers \(t,r\) are amplitudes; the
intensity fractions are \(T=t^2,R=r^2\).
The unitary matrix preserves both photon number
and bosonic commutation relations. Reflection
phases may differ in another convention, with
identical observable predictions.

\begin{center}
\begin{tikzpicture}[>=Latex,scale=.82]
\draw[thick,fill=accent!12] (2.5,.8) rectangle (3.3,2.8);
\draw[->,thick] (.3,2.25)--(2.4,2.25);
\draw[->,thick] (.3,1.35)--(2.4,1.35);
\draw[->,thick] (3.4,2.25)--(5.7,2.25);
\draw[->,thick] (3.4,1.35)--(5.7,1.35);
\node[left] at (.3,2.25){\(a\)};
\node[left] at (.3,1.35){\(b\)};
\node[right] at (5.7,2.25){\(c\)};
\node[right] at (5.7,1.35){\(d\)};
\node at (2.9,.3){unitary mixing};
\end{tikzpicture}
\end{center}

\subsection{Multiphoton calculations use creation operators}
When photons are input in known number states,
write the input with raising operators. For
the balanced convention
\[
a^\dagger=\frac{c^\dagger+d^\dagger}{\sqrt2},
\qquad
b^\dagger=\frac{c^\dagger-d^\dagger}{\sqrt2},
\]
the transformation of \(\ket{1,1}\) is
\[
\begin{aligned}
a^\dagger b^\dagger\ket0
&=\tfrac12(c^\dagger+d^\dagger)
(c^\dagger-d^\dagger)\ket0\\
&=\tfrac12[(c^\dagger)^2-(d^\dagger)^2]\ket0\\
&=\boxed{\frac{\ket{2,0}-\ket{0,2}}{\sqrt2}}.
\end{aligned}
\]
The absence of a \(\ket{1,1}\) term means
zero probability of detecting one photon
in each output. This is Hong--Ou--Mandel
(HOM) interference.

\subsection{Where the HOM cancellation comes from}
Two indistinguishable possibilities lead to
a coincidence: both photons transmit, or
both reflect. Their probability \emph{amplitudes}
have opposite signs and cancel at a balanced splitter.
If the photons have distinguishable wavepackets,
those alternatives no longer interfere completely.
For ideal overlap parameter \(\eta\in[0,1]\),
defined as the squared temporal/spectral-mode overlap,
\[
P(\text{coincidence})=\frac{1-\eta}{2}.
\]
At \(\eta=0\), the classical independent-path
coincidence probability is \(1/2\). At \(\eta=1\),
HOM bunching gives zero.

\begin{worked}{Two photons enter the same port}
For \(\ket{2,0}=(a^\dagger)^2\ket0/\sqrt2\),
\[
\ket{2,0}\longrightarrow
\tfrac12\ket{2,0}+\tfrac1{\sqrt2}\ket{1,1}
+\tfrac12\ket{0,2}.
\]
The output probabilities are \(1/4,1/2,1/4\).
For \(\ket{1,1}\) input, by contrast, the
coincidence probability is \emph{zero}.
Different input states lead to different
interference and photon-counting patterns.
\end{worked}

\subsection{Arbitrary \(n\) photons and coherent input}
A fixed \(n\)-photon input gives the binomial result
\[
\ket{n,0}\mapsto
\sum_{k=0}^n\sqrt{\binom nk}\,t^kr^{n-k}\ket{k,n-k}.
\]
A coherent input \(\ket\alpha_a\ket0_b\)
instead produces a \emph{product} of coherent
outputs \(\ket{t\alpha}_c\ket{r\alpha}_d\)
(up to phases). This follows by substituting
the creation-operator transformation into
the exponential Fock expansion.

\subsection{What a detector measures}
An ideal number-resolving detector distinguishes
\(n=0,1,2,\ldots\). A threshold detector
only reports vacuum or at least one photon.
For \(n\) incident photons and single-photon
detection efficiency \(\eta_d\), the probability
of \emph{no click} is
\[
P(0\ \text{clicks}\mid n)=(1-\eta_d)^n,
\]
neglecting dark counts. Real detection
requires careful treatment of loss, dark counts
and finite time resolution.

\subsection{Homodyne detection accesses field quadratures}
Interfere a weak signal \(a\) with a strong coherent
local oscillator \(\beta=|\beta|e^{i\theta}\).
Subtract the output photodetector currents.
The dominant difference term is proportional to
\[
\beta^*a+\beta a^\dagger
=|\beta|[ae^{-i\theta}+a^\dagger e^{i\theta}].
\]
A local-oscillator phase change selects
a different quadrature rather than a different
photon-number observable.

\begin{checkpoint}
For optical circuits: define input modes,
choose a phase convention, transform operators,
restore the \(\sqrt{n!}\) factors, add
amplitudes, and only then square them.
\end{checkpoint}
